% Job posting: https://ca.linkedin.com/jobs/view/senior-software-programmer-audio-engine-at-epic-games-4410884332
\documentclass[11pt]{article}
\usepackage[left=1in,top=1in,right=1in,bottom=1in]{geometry}
\usepackage[hidelinks]{hyperref}
\usepackage{setspace}
\usepackage[T1]{fontenc}
\usepackage{newtxtext,newtxmath}
\ifdefined\pdfgentounicode
  \input{glyphtounicode}
  \pdfgentounicode=1
\fi
\setstretch{1.1}
\pagenumbering{gobble}
\begin{document}
\pagestyle{empty}

\begin{flushright}
Nishal Stanislaus Silva, PhD \\
Toronto, ON \\
+1 613 200 2030 \\
nishal.silva@hotmail.com \\
\href{https://nishal.xyz}{nishal.xyz} \,|\, \href{https://www.linkedin.com/in/nishal-silva}{linkedin.com/in/nishal-silva} \,|\, \href{https://github.com/ns2max}{github.com/ns2max}
\end{flushright}

\noindent Dear Hiring Manager,

\vspace{0.5em}
\noindent I build real-time audio systems in C++ that ship under hard latency budgets, and that is exactly the discipline your audio engine team needs. During my postdoctoral research at the University of Trento, I owned the full pipeline for a real-time audio pattern detector deployed on Elk Audio OS with VST hosting and OSC control: 14 ms inference on a Raspberry Pi 4, F1 0.76, and 74x faster than the DTW baseline it replaced (published in the Journal of the Audio Engineering Society, 2026). I maintain Nebula, an open-source C++17 audio-feature-extraction library with per-feature latency benchmarks on ARM and Raspberry Pi hardware, and I have hands-on experience with JUCE, VST and OSC as audio-engine tooling. This is the same engineering problem Unreal Engine's audio engine solves at scale: correct, low-latency DSP running inside a hard real-time budget, in production C++.

\vspace{0.5em}
\noindent I bring a second, less common qualification to this work: I have played and performed as a guitarist for over 15 years, holding a Trinity College London Rock \& Pop Guitar, Grade 8. That gives me a trained ear for the artifacts that separate a technically correct audio engine from one that sounds and feels right, a distinction that matters as much in a game engine's mix as it does in a concert hall. I have also built and shipped real-time audio-plus-spatial synchronization systems directly: at Trento, I co-designed a multimodal pipeline unifying audio, BCI/EEG and mixed-reality head/hand tracking across four performers simultaneously, and ran two live concerts on that pipeline with ten Meta Quest 3 headsets in real time. That is real-time audio synchronized against spatial and visual state under live performance conditions, the same coordination problem a game engine's audio-visual sync demands.

\vspace{0.5em}
\noindent I have not shipped inside a game engine before, and I want to be direct about that rather than talk around it. What I bring instead is a research-grade depth in real-time audio DSP under hard constraints, a working C++ audio-engine toolset (JUCE, VST, Elk Audio OS), and a performer's sensibility for audio quality that most engine programmers arrive without. This is a deliberate pivot into games, not a stopgap, and I think that combination of research rigor and musical judgment is a genuinely useful angle for a role building Unreal Engine's audio engine.

\vspace{0.5em}
\noindent I am a Canadian Permanent Resident, eligible to work in Canada without sponsorship, based in Toronto and open to relocating to Vancouver for this role. I am available immediately and would welcome the chance to discuss how my real-time audio engineering background can contribute to your team.

\vspace{1em}
\noindent Sincerely, \\
Nishal Stanislaus Silva

\end{document}
